\documentclass[10pt,letterpaper]{article}
\usepackage[utf8]{inputenc}
\usepackage[intlimits]{amsmath}
\usepackage{amsfonts}
\usepackage{amssymb}
\usepackage{ragged2e}
\usepackage[letterpaper, margin=0.5in]{geometry}
\usepackage{graphicx}
\usepackage{cancel}
\usepackage{mathtools}
\usepackage{tabularx}
\usepackage{arydshln}
\usepackage{tensor}
\usepackage{array}
\usepackage{xcolor}
\usepackage[boxed]{algorithm}
\usepackage[noend]{algpseudocode}
\usepackage{listings}
\usepackage{textcomp}
% \usepackage[pdf,tmpdir,singlefile]{graphviz}
\usepackage{mathrsfs}
\usepackage{bbm}
\usepackage{tikz}
\usepackage{tikz-cd}
\usepackage{enumitem}
\usepackage{arydshln}
\usepackage{relsize}
\usepackage{multirow,multicol}
\usepackage{scalerel}
\usepackage{upgreek}
\usepackage{ifthen}
\usepackage{yhmath}

\usetikzlibrary{bayesnet}
\usetikzlibrary{patterns}
\usetikzlibrary{calc}
\setlist{noitemsep}

%%%%%%%%%%%%%%%%%%%%%%%%%%%%%
% Formatting commands
%%%%%%%%%%%%%%%%%%%%%%%%%%%%%
\newcommand{\n}{\hfill\break}
\newcommand{\nn}{\vspace{0.5\baselineskip}\n}
\newcommand{\up}{\vspace{-\baselineskip}}
\newcommand{\hangblock}[2]{\par\noindent\settowidth{\hangindent}{\textbf{#1: }}\textbf{#1: }\nolinebreak#2}
\newcommand{\lemma}[1]{\hangblock{Lemma}{#1}}
\newcommand{\defn}[1]{\hangblock{Defn}{#1}}
\newcommand{\thm}[1]{\hangblock{Thm}{#1}}
\newcommand{\cor}[1]{\hangblock{Cor}{#1}}
\newcommand{\prop}[1]{\hangblock{Prop}{#1}}
\newcommand{\ex}[1]{\hangblock{Ex}{#1}}
\newcommand{\exer}[1]{\hangblock{Exer}{#1}}
\newcommand{\fact}[1]{\hangblock{Fact}{#1}}
\newcommand{\remark}[1]{\hangblock{Remark}{#1}}
\newcommand{\proven}{\;$\square$\n}
\newcommand{\problem}[1]{\par\noindent{\nolinebreak#1}\n}
\newcommand{\problempart}[2]{\par\noindent\indent{}\settowidth{\hangindent}{\textbf{(#1)} \indent{}}\textbf{(#1) }\nolinebreak#2\n}
\newcommand{\ptxt}[1]{\textrm{\textnormal{#1}}}
\newcommand{\inlineeq}[1]{\centerline{$\displaystyle #1$}}
\newcommand{\pageline}{\up\par\noindent\rule{\textwidth}{0.1pt}}

%%%%%%%%%%%%%%%%%%%%%%%%%%%%%
% Math commands
%%%%%%%%%%%%%%%%%%%%%%%%%%%%%
% Set Theory
\newcommand{\card}[1]{\left|#1\right|}
\newcommand{\set}[1]{\left\{#1\right\}}
\newcommand{\setmid}{\;\middle|\;}
\newcommand{\ps}[1]{\mathcal{P}\left(#1\right)}
\newcommand{\pfinite}[1]{\mathcal{P}^{\ptxt{finite}}\left(#1\right)}
\newcommand{\naturals}{\mathbb{N}}
\newcommand{\N}{\naturals}
\newcommand{\integers}{\mathbb{Z}}
\newcommand{\Z}{\integers}
\newcommand{\rationals}{\mathbb{Q}}
\newcommand{\Q}{\rationals}
\newcommand{\reals}{\mathbb{R}}
\newcommand{\R}{\reals}
\newcommand{\complex}{\mathbb{C}}
\newcommand{\C}{\complex}
\newcommand{\halfPlane}{\mathbb{H}}
\let\HSym\H
\let\H\relax
\newcommand{\H}{\halfPlane}
\newcommand{\comp}{^{\complement}}
\DeclareMathOperator{\Hom}{Hom}
\newcommand{\Ind}{\mathbbm{1}}
\newcommand{\cut}{\setminus}
\DeclareMathOperator{\elem}{elem}

% Differentiable Manifolds
\newcommand{\RP}{\mathbb{RP}}
\newcommand{\CP}{\mathbb{RP}}
\newcommand{\osubset}{\overset{\mathclap{\scalebox{0.5}{\ptxt{open}}}}{\subset}}
\newcommand{\osubseteq}{\overset{\mathclap{\scalebox{0.5}{\ptxt{open}}}}{\subseteq}}
\newcommand{\osupset}{\overset{\mathclap{\scalebox{0.5}{\ptxt{open}}}}{\supset}}
\newcommand{\osupseteq}{\overset{\mathclap{\scalebox{0.5}{\ptxt{open}}}}{\supseteq}}
\newcommand{\pdat}[3]{\left.\pd{#1}{#2}\right|_{#3}}
\DeclareMathOperator{\so}{so}
\DeclareMathOperator{\codim}{codim}
\DeclareMathOperator{\Diff}{Diff}
\let\dSym\d
\let\d\relax
\newcommand{\d}{\partial}
\DeclareMathOperator{\gl}{gl}
\DeclareMathOperator{\Ad}{Ad}
\newcommand{\comm}[1]{\left[#1\right]}
\DeclareMathOperator{\Fix}{Fix}
\DeclareMathOperator{\Gr}{Gr}
\DeclareMathOperator{\Ric}{Ric}
\let\OSymbol\O
\let\O\relax
\DeclareMathOperator{\O}{O}
\newcommand{\overarc}[1]{\wideparen{#1}}
\DeclareMathOperator{\II}{\mathit{I\!\!I}}
\DeclareMathOperator{\Vol}{Vol}
\DeclareMathOperator{\Pf}{Pf}
\DeclareMathOperator{\Laplacian}{\Delta}

% Graph Theory
\let\deg\relax
\DeclareMathOperator{\deg}{deg}
\newcommand{\degp}{\ptxt{deg}^{+}}
\newcommand{\degn}{\ptxt{deg}^{-}}
\newcommand{\precdot}{\mathrel{\ooalign{$\prec$\cr\hidewidth\hbox{$\cdot\mkern0.5mu$}\cr}}}
\newcommand{\succdot}{\mathrel{\ooalign{$\cdot\mkern0.5mu$\cr\hidewidth\hbox{$\succ$}\cr\phantom{$\succ$}}}}
\DeclareMathOperator{\cl}{cl}
\DeclareMathOperator{\affdim}{affdim}

% Probability
\newcommand{\parSymbol}{\P}
\newcommand{\Prob}{\mathbb{P}}
\renewcommand{\P}{\Prob}
\newcommand{\Avg}{\mathbb{E}}
\newcommand{\E}{\Avg}
\DeclareMathOperator{\Var}{Var}
\DeclareMathOperator{\Cov}{Cov}
\DeclareMathOperator{\cov}{cov}
\DeclareMathOperator{\Corr}{Corr}
\DeclareMathOperator{\Unif}{Unif}
\DeclareMathOperator{\Binom}{Binom}
\newcommand{\CI}{\mathrel{\text{\scalebox{1.07}{$\perp\mkern-10mu\perp$}}}}
\DeclareMathOperator{\Ber}{Ber}
\DeclareMathOperator{\Bin}{Bin}
\DeclareMathOperator{\Geom}{Geom}
\DeclareMathOperator{\Poisson}{Poisson}
\DeclareMathOperator{\Exp}{Exp}
\DeclareMathOperator{\Beta}{Beta}
\DeclareMathOperator{\Normal}{N}
\DeclareMathOperator{\DistGamma}{Gamma}
\DeclareMathOperator{\eff}{eff}
\DeclareMathOperator{\AVar}{AVar}

% Standard Math
\newcommand{\inv}{^{-1}}
\newcommand{\abs}[1]{\left|#1\right|}
\newcommand{\ceil}[1]{\left\lceil{}#1\right\rceil{}}
\newcommand{\floor}[1]{\left\lfloor{}#1\right\rfloor{}}
\newcommand{\conj}[1]{\overline{#1}}
\newcommand{\of}{\circ}
\newcommand{\tri}{\triangle}
\newcommand{\inj}{\hookrightarrow}
\newcommand{\surj}{\twoheadrightarrow}
\newcommand{\ndiv}{\nmid}
\renewcommand{\epsilon}{\varepsilon}
\newcommand{\divides}{\mid}
\newcommand{\ndivides}{\nmid}
\DeclareMathOperator{\lcm}{lcm}
\DeclareMathOperator{\sgn}{sgn}
\newcommand{\map}[4]{{\arraycolsep=2pt\begin{array}[t]{rcl}#1 & \to & #2\\ {}#3 & \mapsto & #4\end{array}}}
\newcommand{\bigsum}[2]{\smashoperator[lr]{\sum_{\scalebox{#1}{$#2$}}}}
\DeclareMathOperator{\gcf}{gcf}
\newcommand{\restr}[1]{\left.#1\right|}

% Linear Algebra
\newcommand{\Id}{\textrm{\textnormal{Id}}}
\newcommand{\im}{\textrm{\textnormal{im}}}
\newcommand{\norm}[1]{\abs{\abs{#1}}}
\newcommand{\tpose}{^{T}\!}
\newcommand{\iprod}[1]{\left<#1\right>}
\newcommand{\giprod}{\iprod{\;\,,\;}}
\DeclareMathOperator{\tr}{tr}
\DeclareMathOperator{\trace}{tr}
\newcommand{\chgBasMat}[3]{\!\!\tensor*[_{#1}]{\left[#2\right]}{_{#3}}}
\newcommand{\vecBas}[2]{\tensor*[]{\left[#1\right]}{_{#2}}}
\DeclareMathOperator{\GL}{GL}
\DeclareMathOperator{\Mat}{Mat}
\DeclareMathOperator{\vspan}{span}
\DeclareMathOperator{\rank}{rank}
\newcommand{\V}[1]{\vec{#1}}
\DeclareMathOperator{\proj}{proj}
\DeclareMathOperator{\compProj}{comp}
\DeclareMathOperator{\row}{row}
\newcommand{\smallPMatrix}[1]{\paren{\begin{smallmatrix}#1\end{smallmatrix}}}
\newcommand{\smallBMatrix}[1]{\brack{\begin{smallmatrix}#1\end{smallmatrix}}}
\newcommand{\pmat}[1]{\begin{pmatrix}#1\end{pmatrix}}
\newcommand{\bmat}[1]{\begin{bmatrix}#1\end{bmatrix}}
\newcommand{\dual}{^{*}}
\newcommand{\pinv}{^{\dagger}}
\newcommand{\horizontalMatrixLine}{\ptxt{\rotatebox[origin=c]{-90}{$|$}}}
\DeclareMathOperator{\range}{range}
\DeclareMathOperator{\Symm}{Symm}
\DeclareMathOperator{\SU}{SU}
\DeclareMathOperator{\U}{U}

% Multilinear Algebra
\let\Lsym\L
\let\L\relax
\DeclareMathOperator{\L}{\mathscr{L}}
\DeclareMathOperator{\A}{\mathcal{A}}
\DeclareMathOperator{\Alt}{Alt}
\DeclareMathOperator{\Sym}{Sym}
\newcommand{\ot}{\otimes}
\newcommand{\ox}{\otimes}
\DeclareMathOperator{\asc}{asc}
\DeclareMathOperator{\asSet}{set}
\DeclareMathOperator{\sort}{sort}
\DeclareMathOperator{\ringA}{\mathring{A}}
\DeclareMathOperator{\Sh}{Sh}
\DeclareMathOperator{\Bil}{Bil}

% Topology
\newcommand{\closure}[1]{\overline{#1}}
\newcommand{\uball}{\mathcal{U}}
\DeclareMathOperator{\Int}{Int}
\DeclareMathOperator{\Ext}{Ext}
\DeclareMathOperator{\Bd}{Bd}
\DeclareMathOperator{\rInt}{rInt}
\DeclareMathOperator{\ch}{ch}
\DeclareMathOperator{\ah}{ah}
\newcommand{\LargerTau}{\mathlarger{\mathlarger{\mathlarger{\mathlarger{\tau}}}}}
\newcommand{\Tau}{\mathcal{T}}

% Analysis
\DeclareMathOperator{\Graph}{Graph}
\DeclareMathOperator{\epi}{epi}
\DeclareMathOperator{\hypo}{hypo}
\DeclareMathOperator{\supp}{supp}
\newcommand{\lint}[2]{\underset{#1}{\overset{#2}{{\color{black}\underline{{\color{white}\overline{{\color{black}\int}}\color{black}}}}}}}
\newcommand{\uint}[2]{\underset{#1}{\overset{#2}{{\color{white}\underline{{\color{black}\overline{{\color{black}\int}}\color{black}}}}}}}
\newcommand{\alignint}[2]{\underset{#1}{\overset{#2}{{\color{white}\underline{{\color{white}\overline{{\color{black}\int}}\color{black}}}}}}}
\newcommand{\extint}{\ptxt{ext}\int}
\newcommand{\extalignint}[2]{\ptxt{ext}\alignint{#1}{#2}}
\newcommand{\conv}{\ast}
\newcommand{\pd}[2]{\frac{\partial{}#1}{\partial{}#2}}
\newcommand{\del}{\nabla}
\DeclareMathOperator{\grad}{grad}
\DeclareMathOperator{\curl}{curl}
\let\div\relax
\DeclareMathOperator{\div}{div}
\DeclareMathOperator{\vol}{vol}

% Complex Analysis
\let\Re\relax
\DeclareMathOperator{\Re}{Re}
\let\Im\relax
\DeclareMathOperator{\Im}{Im}
\DeclareMathOperator{\Res}{Res}

% Abstract Algebra
\DeclareMathOperator{\ord}{ord}
\newcommand{\generated}[1]{\left<#1\right>}
\newcommand{\cycle}[1]{\smallPMatrix{#1}}
\newcommand{\id}{\ptxt{id}}
\newcommand{\iso}{\cong}
\DeclareMathOperator{\Aut}{Aut}
\DeclareMathOperator{\SL}{SL}
\DeclareMathOperator{\op}{op}
\newcommand{\isom}[4]{\!\!\!\begin{array}[t]{rcl}#1 & \!\!\!\!\overset{\sim}{\to} & \!\!\!\!#2\\ #3 & \!\!\!\!\mapsto & \!\!\!\!#4\end{array}}
\newcommand{\F}{\mathbb{F}}
\newcommand{\acts}{\;\reflectbox{\rotatebox[origin=c]{-90}{$\circlearrowright$}}\;}
\newcommand{\disjunion}{\mathrel{\text{\scalebox{1.07}{$\perp\mkern-10mu\perp$}}}}
\DeclareMathOperator{\SO}{SO}
\DeclareMathOperator{\stab}{stab}
\DeclareMathOperator{\nullity}{nullity}
\DeclareMathOperator{\Perm}{Perm}
\DeclareMathOperator{\nsubgp}{\vartriangleleft}
\DeclareMathOperator{\notnsubgp}{\ntriangleleft}
\newcommand{\presentation}[2]{\left<#1\;\middle|\;#2\right>}
\DeclareMathOperator{\Char}{char}
\DeclareMathOperator{\fchar}{char}
\DeclareMathOperator{\triv}{triv}
\DeclareMathOperator{\reg}{reg}
\DeclareMathOperator{\std}{std}
\DeclareMathOperator{\Func}{Func}
\DeclareMathOperator{\End}{End}

% Convex Optimization
\let\sectionSymbol\S
\let\S\relax
\newcommand{\S}{\mathbb{S}}
\DeclareMathOperator{\dist}{dist}
\DeclareMathOperator{\dom}{dom}
\DeclareMathOperator{\diag}{diag}
\DeclareMathOperator{\ones}{\mathbbm{1}}
\newcommand{\minimizeOver}[3]{\begin{array}{rl}\underset{#1}{\ptxt{minimize}} & #2\\ \ptxt{subject to} & #3\end{array}}
\newcommand{\maximizeOver}[3]{\begin{array}{rl}\underset{#1}{\ptxt{maximize}} & #2\\ \ptxt{subject to} & #3\end{array}}
\newcommand{\minimizationProblem}[2]{\minimizeOver{}{#1}{#2}}
\newcommand{\maximizationProblem}[2]{\maximizeOver{}{#1}{#2}}
\newcommand{\minimizeOverUnconstrained}[2]{\begin{array}{rl}\underset{#1}{\ptxt{minimize}} & #2\end{array}}
\newcommand{\maximizeOverUnconstrained}[2]{\begin{array}{rl}\underset{#1}{\ptxt{maximize}} & #2\end{array}}
\newcommand{\minimizationUnconstrained}[1]{\minimizeOverUnconstrained{}{#1}}
\newcommand{\maximizationUnconstrained}[1]{\maximizeOverUnconstrained{}{#1}}
\DeclareMathOperator{\argmin}{argmin}
\DeclareMathOperator{\argmax}{argmax}

% Proofs
\newcommand{\st}{s.t.}
\newcommand{\unique}{!}
\newcommand{\iffdef}{\overset{\ptxt{def}}{\Leftrightarrow}}
\newcommand{\eqVertical}{\rotatebox[origin=c]{90}{=}}
\newcommand{\mapsfrom}{\mathrel{\reflectbox{\ensuremath{\mapsto}}}}
\newcommand{\mapsdown}{\rotatebox[origin=c]{-90}{$\mapsto$}\mkern2mu}
\newcommand{\mapsup}{\rotatebox[origin=c]{90}{$\mapsto$}\mkern2mu}
\newcommand{\from}{\!\mathrel{\reflectbox{\ensuremath{\to}}}}
\newcommand{\labeledeq}[1]{\overset{\mathclap{\ptxt{#1}}}{=}}
\newcommand{\eqdef}{\labeledeq{def}}
\newcommand{\leftsquigarrow}{\mathrel{\reflectbox{\ensuremath{\rightsquigarrow}}}}
\newcommand{\selfmap}{\mathbin{\scalebox{.85}{\lefteqn{\scalebox{.5}{$\blacktriangleleft$}}\raisebox{.34ex}{$\supset$}}}}

% Brackets
\newcommand{\paren}[1]{\left(#1\right)}
\renewcommand{\brack}[1]{\left[#1\right]}
\renewcommand{\brace}[1]{\left\{#1\right\}}
\newcommand{\ang}[1]{\left<#1\right>}

% Algorithms
\algrenewcommand{\algorithmiccomment}[1]{\hskip 1em \texttt{// #1}}
\algrenewcommand\algorithmicrequire{\textbf{Input:}}
\algrenewcommand\algorithmicensure{\textbf{Output:}}
\newcommand{\algP}{\ptxt{\textbf{P}}}
\newcommand{\algNP}{\ptxt{\textbf{NP}}}
\newcommand{\algNPC}{\ptxt{\textbf{NP-Complete}}}
\newcommand{\algNPH}{\ptxt{\textbf{NP-Hard}}}
\newcommand{\algEXP}{\ptxt{\textbf{EXP}}}
\DeclareMathOperator{\fl}{fl}

%%%%%%%%%%%%%%%%%%%%%%%%%%%%%
% Other commands and shorthand
%%%%%%%%%%%%%%%%%%%%%%%%%%%%%
\newcommand{\flag}[1]{\textbf{\textcolor{red}{#1}}}
\let\uSym\u
\let\u\relax
\newcommand{\u}[1]{\underline{#1}}
\let\bSym\b
\let\b\relax
\newcommand{\b}[1]{\textbf{#1}}
\let\iSym\i
\let\i\relax
\newcommand{\i}[1]{\textit{#1}}
\let\scSym\sc
\let\sc\relax
\newcommand{\sc}[1]{\textsc{#1}}
\newcommand{\mf}[1]{\mathfrak{#1}}
\newcommand{\mc}[1]{\mathcal{#1}}
\newcommand{\ms}[1]{\mathscr{#1}}

%%%%%%%%%%%%%%%%%%%%%%%%%%%%%%%%%%%%%%%
% Make l's curvy in math environments %
%%%%%%%%%%%%%%%%%%%%%%%%%%%%%%%%%%%%%%%
\mathcode`l="8000
\begingroup
\makeatletter
\lccode`\~=`\l
\DeclareMathSymbol{\lsb@l}{\mathalpha}{letters}{`l}
\lowercase{\gdef~{\ifnum\the\mathgroup=\m@ne \ell \else \lsb@l \fi}}%
\endgroup

%%%%%%%%%%%%%%%%%%%%%%%%%
% Fix \vdots and \ddots %
%%%%%%%%%%%%%%%%%%%%%%%%%
\usepackage{letltxmacro}
\LetLtxMacro\orgvdots\vdots
\LetLtxMacro\orgddots\ddots

\makeatletter
\DeclareRobustCommand\vdots{%
	\mathpalette\@vdots{}%
}
\newcommand*{\@vdots}[2]{%
	% #1: math style
	% #2: unused
	\sbox0{$#1\cdotp\cdotp\cdotp\m@th$}%
	\sbox2{$#1.\m@th$}%
	\vbox{%
		\dimen@=\wd0 %
		\advance\dimen@ -3\ht2 %
		\kern.5\dimen@
		% remove side bearings
		\dimen@=\wd2 %
		\advance\dimen@ -\ht2 %
		\dimen2=\wd0 %
		\advance\dimen2 -\dimen@
		\vbox to \dimen2{%
			\offinterlineskip
			\copy2 \vfill\copy2 \vfill\copy2 %
		}%
	}%
}
\DeclareRobustCommand\ddots{%
	\mathinner{%
		\mathpalette\@ddots{}%
		\mkern\thinmuskip
	}%
}
\newcommand*{\@ddots}[2]{%
	% #1: math style
	% #2: unused
	\sbox0{$#1\cdotp\cdotp\cdotp\m@th$}%
	\sbox2{$#1.\m@th$}%
	\vbox{%
		\dimen@=\wd0 %
		\advance\dimen@ -3\ht2 %
		\kern.5\dimen@
		% remove side bearings
		\dimen@=\wd2 %
		\advance\dimen@ -\ht2 %
		\dimen2=\wd0 %
		\advance\dimen2 -\dimen@
		\vbox to \dimen2{%
			\offinterlineskip
			\hbox{$#1\mathpunct{.}\m@th$}%
			\vfill
			\hbox{$#1\mathpunct{\kern\wd2}\mathpunct{.}\m@th$}%
			\vfill
			\hbox{$#1\mathpunct{\kern\wd2}\mathpunct{\kern\wd2}\mathpunct{.}\m@th$}%
		}%
	}%
}
\makeatother

\newcommand{\B}{
	\begin{tikzpicture}
	\filldraw [fill=red, draw=black] (0, 0) rectangle (0.37, 0.45);
	\draw [line width=0.5mm, white ] (0.1,0.08) -- (0.1,0.38);
	\draw[line width=0.5mm, white ] (0.1, 0.35) .. controls (0.2, 0.35) and (0.4, 0.2625) .. (0.1, 0.225);
	\draw[line width=0.5mm, white ] (0.1, 0.225) .. controls (0.2, 0.225) and (0.4, 0.1625) .. (0.1, 0.1);
	\end{tikzpicture}
}

% Allow custom symbols as arrows in tikz-cd
\tikzset{
  symbol/.style={
    draw=none,
    every to/.append style={
      edge node={node [sloped, allow upside down, auto=false]{$#1$}}}
  }
}
% Example: \arrow[r,symbol=\cong]
% https://tex.stackexchange.com/questions/394154/how-to-include-inclusion-subgroup-relationship-in-tikz-cd-diagram

\author{Thomas Cohn}
\title{Math 635 Lecture 37}
\date{4/16/21} % Can also use \today

\begin{document}
\maketitle
\setlength\RaggedRightParindent{\parindent}
\RaggedRight

\par\noindent
Recall from last time: Let $M$ be a compact, oriented, Riemannian manifold. Then the Laplacian on forms $\Laplacian:\Omega^{k}\to\Omega^{k}$ is $\Laplacian=\delta{}d+d\delta$. (Note that these are technically different $d$'s and different $\delta$'s, because $d$ increases the degree and $\delta$ reduces it.) And we have the deRham complex
\[
	\begin{tikzcd}
		\cdots \arrow[r, shift left, "d"] & \Omega^{k-1} \arrow[l, shift left, "\delta"] \arrow[r, shift left, "d"] & \Omega^{k} \arrow[l, shift left, "\delta"] \arrow[r, shift left, "d"] & \Omega^{k+1} \arrow[l, shift left, "\delta"] \arrow[r, shift left, "d"] & \cdots \arrow[l, shift left, "\delta"]
	\end{tikzcd}
\]
where $d^{*}=\delta$, i.e., $\forall\alpha\in\Omega^{k-1},\beta\in\Omega^{k}$, $\iprod{d\alpha,\beta}_{l^{2}}=\iprod{\alpha,\delta\beta}_{l^{2}}$.\n

\lemma{
	\begin{enumerate}[leftmargin=4\parindent, label=(\arabic*)]
		\item $\delta^{2}=0$ (because $\delta^{2}=(d^{2})^{*}$)
		\item $\Laplacian^{*}=\Laplacian$: $\forall\alpha,\beta\in\Omega^{k}$, $\iprod{\Laplacian\alpha,\beta}_{l^{2}}=\iprod{\alpha,\Laplacian\beta}_{l^{2}}$
		\item $\comm{\Laplacian,d}=0$, $\comm{\Laplacian,\delta}=0$
		\item $\Laplacian\alpha=0$ iff $d\alpha=0$ and $\delta\alpha=0$
	\end{enumerate}\up\nn
	Proof (3): $(\delta{}d+d\delta)d=\cancel{\delta{}dd}+d\delta{}d$ and $d(\delta{}d+d\delta)=d\delta{}d+\cancel{dd\delta}=d\delta{}d$. So $\comm{\Laplacian,d}=0$. (Identical proof for $\comm{\Laplacian,\delta}=0$.)\proven
	(4): $\Leftarrow$ is obvious. For $\Rightarrow$, note that $\Laplacian\alpha=0$ implies
	\[
		0=\iprod{\Laplacian\alpha,\alpha}=\iprod{\delta{}d\alpha,\alpha}+\iprod{d\delta\alpha,\alpha}=\iprod{d\alpha,d\alpha}+\iprod{\delta\alpha,\delta\alpha}=\norm{d\alpha}_{l^{2}}^{2}+\norm{\delta\alpha}_{l^{2}}^{2}
	\]
	So $d\alpha=0$ and $\delta\alpha=0$.\proven
}

\subsection*{General Things about Linear Differential Operators}

\par\noindent
Let $U\subseteq\R^{n}$, consider $C^{\infty}(U,\C)$.\n

\defn{
	A \u{differential operator} $P$ on $C^{\infty}(U,\C)$ is of the form $\forall{}f\in{}C^{\infty}(U,\C)$,
	\[
		P(f)=\sum_{\alpha}c_{\alpha}(x)(D^{\alpha}f)(x)\quad\alpha\;\ptxt{multi-index with}\;\abs{\alpha}\le{}n\qquad{}D^{\alpha}=\frac{1}{i}\paren{\pd{}{x_{1}}}^{\alpha_{1}}\cdots\frac{1}{i}\paren{\pd{}{x_{n}}}^{\alpha_{n}}
	\]
	(here, we're using $i=\sqrt{-1}$). Note that $P$ is local -- $P(f)(x)$ only depends on $f$ in a neighborhood of $x$.\n
}

\par\noindent
On manifolds, on $C^{\infty}(M,\C)$, linear differential operators are generated as a ring by multiplying by functions. For vector bundles/systems, we have $P:C^{\infty}(U,\C^{r})\to{}C^{\infty}(U,\C^{s})$, where $P=(P_{ij})_{s\times{}r}$, and $P_{ij}$ is a scalar differential operator.\n

\par\noindent
In the manifold setting, suppose we have two bundles $\mc{E}$ and $\mc{F}$:
\[
	\begin{tikzcd}
		\mc{E} \arrow[dr] &[-20pt] &[-20pt] \mc{F} \arrow[dl]\\
		 & M
	\end{tikzcd}
\]
Consider $P:\Gamma(\mc{E})\to\Gamma(\mc{F})$. Under local trivializations on the same $U\subseteq{}M$, we have $\Gamma(\restr{\mc{E}}_{U})\cong{}C^{\infty}(U,\C^{r})\cong\Gamma(\restr{\mc{F}}_{U})$. So $P$ is locally the Euclidean case.\n

\par\noindent
We already know a bunch of examples!\n

\ex{
	$d:\Omega^{k}\to\Omega^{k+1}$ and $\nabla:\Gamma(\mc{E})\to\Gamma(\mc{E}\otimes{}T^{*}M)$ are differential operators of order $1$.\n
}

\par\noindent
The main invariant associated to a differential operator $P$ is called its symbol. The symbol captures the top degree part of the operator. For an operator on $C^{\infty}(U,\C)$, think of computing: take $x_{0}\in{}U,\xi_{0}\in{}T_{x_{0}}^{*}U$ (i.e. $(x_{0},\xi_{0})\in{}T^{*}U$). Pick $\chi,\rho$ functions, both $C^{\infty}_{0}(U,\R)$ such that $\chi\equiv{}1$ near $x_{0}$ and $d\rho_{x_{0}}=\xi_{0}$. Consider $P(\chi{}e^{i\tau\rho})(x_{0})$ for $\tau\gg1$. We get
\[
	P(\chi{}e^{i\tau\rho})(x_{0})=\sum_{\mathclap{\abs{\alpha}=m\eqdef\deg{}P}}c_{\alpha}(x_{0})D^{\alpha}(\chi{}e^{i\tau\rho})+\ptxt{lower degree terms}
\]
Look for the highest power of $\tau$ on the right-hand side.
\[
	\tau^{m}e^{i\tau\rho(x_{0})}\sum_{\abs{\alpha}=m}c_{\alpha}(x_{0})\underbrace{(\del\rho(x_{0}))^{\alpha}}_{\xi_{1}^{\alpha_{1}}\cdots\xi_{n}^{\alpha_{n}}}+\ptxt{lower order terms in $\tau$}
\]
Now, forget the $\tau$ and forget the exponential.\n

\defn{
	The \u{symbol} of a differential operator $P$ is
	\[
		\sigma_{P}(x_{0},\xi_{0})=\sum_{\abs{\alpha}=m}c_{\alpha}(x_{0})\xi^{\alpha}=\lim_{\tau\to\infty}e^{-i\tau\rho(x_{0})}\frac{1}{\tau^{m}}P(\chi{}e^{i\tau\rho})(x_{0})
	\]
}

\par\noindent
Conclusion: On manifolds, $P:C^{\infty}(M,\C)\selfmap$, $\sigma_{P}:T^{*}M\to\C$ is well-defined. Specifically, it's a homogeneous polynomial in $\chi$ of degree $M$ on each fiber.\n

\ex{
	Let $X\in\mf{X}(M)$. Then $\sigma_{\mc{L}_{X}}:T^{*}M\to\C$ is $\sigma_{\mc{L}_{X}}(p,\xi)=i\iprod{\xi,X_{p}}$.\n
}

\ex{
	$\Laplacian:C^{\infty}(M)\selfmap$, $\sigma_{\Laplacian}:T^{*}M\to\C$ is $\sigma_{\Laplacian}(p,\xi)=g^{ij}(p)\xi_{i}\xi_{j}=\norm{\xi}^{2}\Id$.\n
	(Recall: this is from the definition of $\Laplacian$ in coordinates.)\n
}

\par\noindent
If $P:\Gamma(\mc{E})\to\Gamma(\mc{F})$,
\[
	\begin{tikzcd}
		\mc{E} \arrow[dr] &[-20pt] &[-20pt] \mc{F} \arrow[dl]\\
		 & M
	\end{tikzcd}
\]
Then for $(p,\xi)\in{}T^{*}M$, $\sigma_{P}(p,\xi):\mc{E}_{p}\to\mc{F}_{p}$ is a linear map between the fibers.\n

\prop{
	If $P$ and $Q$ are differential operators such that $P\of{}Q$ makes sense, then $\sigma_{P\of{}Q}=\sigma_{P}\of\sigma_{Q}$.\n
}

\defn{
	$P$ is an \u{elliptic operator} iff $\forall{}(p,\xi)$ with $\xi\ne{}0$, $\sigma_{P}(p,\xi)$ is invertible.\n
}

\ex{
	$P=\Laplacian$ is elliptic. $\sigma_{\Laplacian}=\norm{\xi}^{2}\Id$, so it's invertible everywhere. $\Laplacian$ has an approximate inverse $G$ -- ``$\sigma_{G}(p,\xi)=\frac{\Id}{\norm{\xi}^{2}}$''. $\Laplacian\of{}G-I$ and $G\of\Laplacian-I$ are smoothing operators -- they're very small.\n
}

\par\noindent
We're skipping a lot of stuff here, but...\n

\thm{
	(Spectral Theorem of the Laplacian) Consider $\Laplacian:\Omega^{k}(M)\to\Omega^{k}(M)$, where $M$ is a compact, oriented, Riemannian manifold. Then
	\begin{enumerate}[label=(\arabic*), leftmargin=4\parindent]
		\item $\mc{U}^{k}=\ker\Laplacian$ has finite dimension.
		\item There's an orthonormal basis (in the $l^{2}$ sense) of $\Omega^{k}$, $\set{\alpha_{j}}$, and $0\le\lambda_{0}\le\lambda_{1}\le\cdots\to+\infty$ \st{} $\Laplacian\alpha_{j}=\lambda_{j}\alpha_{j}$. In other words, we can think of $\Laplacian$ as an infinite matrix, that can be diagonalized by $\alpha_{j}$ and $\lambda_{j}$. We'll write $\forall\alpha\in\Omega^{k}$, $\alpha=\alpha^{H}+\sum_{\lambda>0\;\ptxt{distinct}}\alpha_{\lambda}$. $\alpha^{H}$ is the harmonic piece -- $\Laplacian\alpha^{H}=0$, $\Laplacian\alpha_{\lambda}=\lambda\alpha_{\lambda}$. We define Green's operator $G(\alpha)=\sum_{\lambda>0}\frac{1}{\lambda}\alpha_{\lambda}$. So $(I-\Laplacian\of{}G)(\alpha)=G^{H}$.
	\end{enumerate}
}

\ex{
	For $\Laplacian$ on $C^{\infty}(S^{2})$, we have $\lambda=k(k+1)$ with multiplicity $2k+1$ (for $k\in\Z_{\ge{}0}$).\n
}

\end{document}